\section{System and Variants}\label{sec:method}

\textbf{Base system.} The assistant as it stood before this study uses no language model to decide. A question is split into segments (one
per requested procedure), each normalized (abbreviations, synonyms, the fields asked for such as documents or fees), and procedures are
ranked by inverse document frequency over whole Vietnamese syllables\cite{robertson2009}, with a penalty for mismatched diacritics, a bonus
for adjacent matches, and thresholds on score and coverage (the share of the question's weight found in the name). A policy module checks
the plan against the data (real identifiers, commune scope, province versions, at most three procedures) and outputs the behavior. Answers
are composed by code from the requested fields of the record, with the source. An optional LLM rewriting step is off in all experiments:
we measure the decision, not the wording.

\textbf{Rule layer (A).} Rules written from the base system's errors on the development set, each behind a switch; with all switches off,
the system reproduces the base output on every test question.
\begin{itemize}
\item \emph{Letter stripping}: remove greetings, the addressed agency, self-introductions and thanks, which otherwise match procedure names.
\item \emph{Citizen lexicon}: about twenty patterns map citizen wording to current official names, including terms changed by recent laws
(the household registration book, abolished by the 2020 Law on Residence; ``citizen identity card'', renamed ``identity card'' in 2023).
\item \emph{Weak-name filter}: a segment sharing a few words with a very long procedure name does not name it.
\item \emph{Out-of-scope gate}: patterns for the status of one's own file, complaints and petitions, law and penalties, police internal
matters and recruitment, and topics absent from the data yield \emph{refuse}; explicit procedural questions (how to file a complaint) are
exempt.
\item \emph{Focus}: in long letters, keep only segments that ask something.
\item \emph{Organization default}: without a mentioned organization, prefer the household or individual version.
\end{itemize}

\textbf{Second round (A2).} Once the first test set had been run and its results read, a second round, behind one more switch, was written
on the remaining errors of A on the development and first test sets. It makes the lexicon tolerate inserted words; lets a request
recognized by the lexicon override a procedure inferred from a life event in the story (a parent's death suggests death registration, but
the question asks to transfer the parent's car); adds out-of-scope patterns (status in plain words, police staff asking about their benefits,
private disputes); answers questions on the processing time set by regulation;
and answers with the version common to all provinces instead of asking for the province, as the labeling guideline prescribes.

\textbf{Reranker (B).} Candidates are the rule ranker's top eight and the ten in-scope procedures closest to the question under
BGE-M3\cite{chen2024m3}; a cross-encoder, bge-reranker-v2-m3\cite{bgererank}, scores each (question, procedure description) pair. The
system refuses when the best score is below a threshold and replaces the rule's choice when another candidate wins by a margin; alone, it
also answers when the rules find nothing but the best score is positive. Thresholds were chosen by grid search on development (alone:
$-3.5$ and $0.5$; on top of A: $-8$ and $2$, no rescue).

\textbf{LLM reader (C).} A local 4-billion-parameter model (Qwen3-4B-Instruct-2507, Q4\_K\_M quantization, Ollama, temperature 0, JSON
schema)\cite{qwen3,qwen3inst} returns one of six intents (procedure, file status, complaint, law, internal, other) and, for procedure
questions, the procedure's name in words, which the rule ranker maps to a procedure; the model never writes an identifier. Non-procedure
intents become refusals; the \emph{name-only} variant ignores the intent.

\textbf{RAG baseline (R).} To represent a common design in which the model decides, BGE-M3 retrieves the five in-scope procedures closest to
the raw question and the same 4B model returns \emph{answer} with a number, \emph{clarify} or \emph{refuse}; no rule is used. A first, short
prompt asked for clarification too often on development; two variants (ordered decision rules, and the same with procedure descriptions)
were then declared together before either was run. Their development balanced accuracies were 0.343, 0.445 and 0.432, and the second was
run once on each test set. R was added after the first test set had been run, but nothing in it was chosen on a test set.
